\subsection{Parte 2}
\subsubsection{Ejercicio 1}
\begin{exercise}
Dada una cadena de Markov homogénea con espacio de estados $E=\{0,1\}$ y matriz de transición
\[
\mathbb{P}=
\begin{pmatrix}
1-\alpha & \alpha \\
\beta & 1-\beta
\end{pmatrix},
\]
calcula $P_0(T_0=n)$ y $E_0[T_0]$, donde $T_0=\inf\{n\ge1 : X_n=0\}$ es el tiempo de primer retorno al estado $0$.
\end{exercise}

La matriz de transición es irreducible y el espacio de estados $E$ es finito, por lo tanto la cadena es persistente positiva y tendrá una distribución estacionaria.

\[
\pi_0=\frac{1}{E_0[T_0]}, \qquad \pi_1=\frac{1}{E_1[T_1]}.
\]

Para hallar $\pi$, usamos las ecuaciones de equilibrio global:
\[
\left\{
\begin{array}{l}
\pi \mathbb{P} = \pi, \\
\pi \mathbb{1}= \pi \longrightarrow \pi_0 + \pi_1 = 1.
\end{array}
\right.
\]

Desarrollando la primera ecuación:
\[
\pi_0 = (1-\alpha)\pi_0 + \beta\pi_1
\quad \Longrightarrow \quad
\alpha\pi_0 = \beta\pi_1
\quad \Longrightarrow \quad
\frac{\pi_0}{\pi_1} = \frac{\beta}{\alpha}.
\]

Usando $\pi_0 + \pi_1 = 1$:
\[
\pi_0 = \frac{\beta}{\alpha+\beta}, \qquad \pi_1 = \frac{\alpha}{\alpha+\beta}.
\]

Como $\pi_0 = \dfrac{1}{E_0[T_0]}$, se tiene:
\[
E_0[T_0] = \frac{1}{\pi_0} = \frac{\alpha+\beta}{\beta}.
\]

\medskip
\textbf{Cálculo de $P_0(T_0=n)$}

Partiendo del estado $0$, la cadena puede retornar a $0$ en $n$ pasos sólo si:
\[
X_1=1, \, X_2=1, \dots, X_{n-1}=1, \, X_n=0.
\]
% Probabilidades de retorno al estado 0 (solo la parte solicitada)
\[
P_0(T_0=1)=\mathbb{P}(X_1=0\mid X_0=0)=p_{00}=1-\alpha.
\]
\[
P_0(T_0=2)=\mathbb{P}(X_1=1,X_2=0\mid X_0=0)=p_{01}p_{10}=\alpha\beta.
\]

Por tanto:
\[
P_0(T_0=n)=
\begin{cases}
1-\alpha, & n=1,\\[4pt]
\alpha(1-\beta)^{\,n-2}\beta, & n\ge2.
\end{cases}
\]

Comprobamos:
\[
E_0[T_0] = \sum_{n=2}^{\infty} n P_0(T_0=n)
= \sum_{n=2}^{\infty} n\, \alpha(1-\beta)^{n-2}\beta
= \frac{\alpha+\beta}{\beta},
\]
lo cual coincide con el valor obtenido a partir de la distribución estacionaria.



\subsubsection{Ejercicio 2}
\subsubsection{Ejercicio 7}

\begin{exercise}
    En ajedrez el rey puede moverse en horizontal, vertical o diagonal, con desplazamientos de una celda sobre un tablero de $8\times 8$. Supongamos que $X_n$ es la posición del rey si en cada paso se elige uno de los movimientos legales al azar.  Calcula la distribución estacionaria de $X_n$ y el tiempo medio de retorno a una esquina cuando se parte de ella.
\end{exercise}

Para calcular la distribución estacionaria partimos de las ecuaciones d eequilibrio global: $\left\{\begin{array}{l}
    \pi=\pi\mathbb{P}\\
    \pi\mathds{1}=1
\end{array}\right.$\\

También tenemos que tener en cuenta que $g_i$ es el grado del nodo $i$. De la primera ecuación de equilibrio sacamos que:

$$
\pi_i\frac{1}{g_i}=\pi_j\frac{1}{g_j}=cte.\Rightarrow\ \frac{\pi_i}{g_i}=cte.\ \forall i\in E
$$

De la segunda ecuación de equilibrio sacamos que:

$$
\sum_{i\in E}\pi_i=1\Rightarrow\sum_{i\in E}cte.\ g_i=1\Rightarrow cte.=\frac{1}{\sum_{i\in E}g_i}
$$

Volviendo a la ecuación anterior, podemos obtener una fórmula para la distribución estacionaria:

$$
\pi_j=\frac{g_j}{\sum_{i\in E}g_i}
$$

A partir de este punto, solo tenemos que calcular el número total de nodos y el grado de cada uno de ellos. Como tenemos un tablero $8\times 8$, 
\begin{itemize}
    \item Tenemos 4 esquinas, cuyo grado es 3.
    \item Tenemos 6+6+6+6=24 casillas laterales, cuyo grado es 5
    \item Tenemos 64-24-4=36 casillas centrales, cuyo grado es 8
\end{itemize}

En total tendremos un grado total de $4\times3+24\times5+36\times8=420$, por lo que la distribución estacionaria será

$$
\pi_i=\frac{g_i}{420}
$$

\subsubsection{Ejercicio 8}
\subsubsection{Ejercicio 10}
\subsubsection{Ejercicio 14}
\subsubsection{Ejercicio 17}

